% Redrawn 2026-10-15; main labels 8 pt, notes 7.5 pt; actual fonts embedded in PDF.
% export-reference-width-mm: 112
\begin{tikzpicture}[font=\figurefont\fontsize{8}{10}\selectfont,text=BookInk,
  port/.style={draw=FigureStroke,fill=FigureBlueFill,line width=.6pt,
    rounded corners=1.8pt,minimum width=24mm,minimum height=8mm,align=center},
  flow/.style={-{Latex[length=2mm]},draw=FigureStroke,line width=.6pt,
    rounded corners=1.8pt}]
  \foreach \offset/\idx in {0/1,5.5/2}{
    \begin{scope}[xshift=\offset cm]
      \fill[FigureGreenFill] (0,0) rectangle (2.4,1.6);
      \draw[BookMuted,line width=.5pt,xstep=.8cm,ystep=.8cm]
        (0,0) grid (2.4,1.6);
      \node at (.4,.4) {$s_{\idx}$};
      \node at (2,1.2) {$g_{\idx}$};
      \draw[flow,BookTeal] (.68,.4) -- (1.2,.4) -- (1.2,1.2) -- (1.73,1.2);
      \node at (1.2,2.02) {子网格$\idx$};
      \node[port] (a\idx) at (1.2,-.8) {动作接口$a_{\idx}$};
      \draw[flow] (a\idx.north) -- (1.2,0);
    \end{scope}
  }
  \node[text=BookMuted,font=\FigureNoteFont\fontsize{7.5}{9.5}\selectfont] at (3.95,2.65)
    {无共享限制：$\mathsf A=\mathsf A_1\times\mathsf A_2$，分别选动作};
  \node[port,draw=BookAmber,fill=FigureYellowFill,minimum width=60mm]
    (schedule) at (3.95,-2.75) {加入共享调度：每步最多激活一个};
  \draw[flow,BookAmber,dashed] (schedule.north) -- (3.95,-1.85);
  \draw[BookAmber,dashed,line width=.6pt] (1.2,-1.85) -- (6.7,-1.85);
  \draw[flow,BookAmber,dashed] (1.2,-1.85) -- (a1.south);
  \draw[flow,BookAmber,dashed] (6.7,-1.85) -- (a2.south);
  \node[font=\FigureNoteFont\fontsize{7.5}{9.5}\selectfont,text=BookMuted,font=\FigureNoteFont\fontsize{7.5}{9.5}\selectfont] at (3.95,-3.55)
    {两接口互斥；不能同时实施各自的局部最优激活动作};
\end{tikzpicture}
